\documentclass[11pt,a4paper]{article}
\usepackage[margin=23mm]{geometry}
\usepackage{amsmath,amssymb,amsthm}
\setlength{\emergencystretch}{2em}
\newtheorem{lemma}{Lemma}
\newtheorem{theorem}[lemma]{Theorem}
\newcommand{\gen}[1]{\langle#1\rangle}
\newcommand{\Alt}{\operatorname{Alt}}
\title{Closing the case of two opposite six-point heads}
\author{Human Compression project: L044 continuation}
\date{}
\begin{document}\maketitle
Retain the complete critical-head hypotheses and their acquired
proofs, including the smaller-degree alternating induction that
produced this configuration. Thus $G\le\Alt(\Omega)$ is transitive
with a connected string C-marking $\rho_0,\ldots,\rho_{r-1}$,
all maximal omissions are intransitive, and $n=2r\ge18$.
Suppose the first and last three marks are actual supported,
extendable six-point Table~2 heads. Write their disjoint supports
as $D_L,D_R$, with exit vertices $v_L,v_R$. The complete
critical-head theorem and its full reversal give
\[
 G_{\ge3}=\Alt(\Omega\setminus(D_L\setminus\{v_L\})),\qquad
 C:=G_{\le r-4}=\Alt(V),\quad
 V=\Omega\setminus(D_R\setminus\{v_R\}).                 \tag{1}
\]
All groups in (1) are actual natural permutation subgroups.
On the left use the literal notation
\[
 \tau=\rho_3=(v_L\ \beta)\delta,\quad \sigma=\rho_4,
 \quad H_C=\gen{\rho_5,\ldots,\rho_{r-4}},\quad
 K_C=\gen{\rho_6,\ldots,\rho_{r-4}}.
\]
The odd involution $\delta$ fixes $D_L,\beta$ and centralizes
$G_{\ge5}$. The first three marks are supported on $D_L$;
$\sigma$ fixes $D_L$; $G_{\ge5}$ fixes $D_L\cup\{\beta\}$.
Every use of eligibility below refers to the full group specified.

\begin{lemma}[The central colours are paired]
Every colour $3,\ldots,r-4$ has two edges eligible for its
full ambient omission in $G$.
\end{lemma}
\begin{proof}
The complete transitive-$B$ exclusion and early-cut calculation
give paired colours $3,4$, and their full reversals give paired
colours $r-4,r-5$. It remains to consider $5\le i\le r-6$.
If colour $i$ had just one eligible edge, the acquired
pure-prefix orbit lemma would make every mark after $i$ fix
the omission orbit containing $D_L$ pointwise. Apply the entire
lemma in reversed order: every mark before $i$ fixes the other
omission orbit, containing $D_R$, pointwise. These are the two
omission orbits, so the cut has actual supported transitive
prefix and suffix and is unshared.

The complete later-unshared-cut theorem applies to this literal
cut. Its only residual indices are $i=r-3$ and $i=r-1$,
both outside the present range. Thus no such unique cut exists.
There is at least one eligible edge at every colour, because
the full group is transitive and its omission is intransitive.
Consequently every colour in the range is paired.
\end{proof}

\begin{lemma}[An odd prefix and the actual attachment]
The complete odd-degree separator arguments apply to $C$ on $V$.
They give an index $s$ with $5\le s<r-4$ such that
\[
             |\operatorname{Fix}_{\Omega}(G_{>s})|\ge2s+1.
                                                               \tag{2}
\]
Their local fixed points survive in this full ambient fixed set.
\end{lemma}
\begin{proof}
The degree and rank of the actual marking of $C$ are
$|V|=n-5=2(r-3)+1$ and $p=r-3\ge6$. Restriction to $V$
is faithful, so actual IP, rank and connectedness are retained.
If a maximal omission $C_i$ were transitive on $V$, adjoining
the right head, whose support meets $V$ at exactly $v_R$,
would make $G_i$ transitive on $\Omega$. Hence all its maximal
omissions are intransitive.

Intersect the two actual natural groups in (1). Full IP gives
\[
 C_{\ge3}=\gen{\rho_3,\ldots,\rho_{r-4}}
     =\Alt(V\setminus(D_L\setminus\{v_L\})).
\]
The literal left head and its odd $\delta$ are unchanged:
every mark of $C$ fixes $\Omega\setminus V$, and the core
transposition is inside $V$. The previous lemma supplies two
ambient $G_i$-eligible edges for every $5\le i\le r-4$.
These edges lie in $V$ and are $C_i$-eligible because
$C_i\le G_i$. This is precisely the valid direction of
eligibility inheritance. All hypotheses of the complete odd
separator are now paid.

It yields $s<p-1=r-4$ and at least $2s+1$ points of $V$
fixed by $\gen{\rho_{s+1},\ldots,\rho_{r-4}}$.
The right exit vertex $v_R$ is not among them: the literal
right crossing $\rho_{r-4}$ moves it. Each of the three
remaining marks is supported on $D_R$ and hence fixes
$V\setminus\{v_R\}$ pointwise. Thus all these local fixed
points are fixed by the full $G_{>s}$, proving (2).
\end{proof}

\begin{lemma}[The exterior involution fixes the later support]
For the separator $s$ just obtained, put
$R=G_{>s}$, $U=\operatorname{Supp}(R)$, $F=\Omega\setminus U$,
and $d=r-s-1$. Then
\[
                         |F|\le2s+2,\qquad \delta|_U=1.
                                                               \tag{3}
\]
\end{lemma}
\begin{proof}
Here $d\ge4$, and the full reversed pure-prefix theorem
applies to $R$. Its orbit $A$ containing $D_R$ has at least
seven points and carries the pure natural $\Alt(A)$.
Let $q\ge1$ be the number of nontrivial $R$-orbits.
The actual marked intersection
\[
 J=R\cap C=\gen{\rho_{s+1},\ldots,\rho_{r-4}}
\]
has rank $d-3$. The exact reversed prefix intersection formula
removes just the five points $D_R\setminus\{v_R\}$ from
the head orbit and its support, leaving every exterior orbit
unchanged. Thus its support has size $|U|-5$ and it has $q$
nontrivial orbits. The reduced head orbit $A'=A\cap V$ is
nontrivial even when $|A|=7$: its two points are exchanged
by the literal $\rho_{r-4}$. For larger $A$ its pure alternating
subgroup is transitive on $A'$.

All $d-3$ colours of $J$ are paired by Lemma~1. The full
ambient interval inequality therefore gives
\[
       2(d-3)\le |U|-5-q,\qquad |U|\ge2d-1+q.          \tag{4}
\]
On the other hand (2) gives $|U|\le n-2s-1=2d+1$.
Consequently $q\le2$, and (4) already proves
$|F|\le2s+2$.

The involution $\delta$ centralizes $R$ and hence its pure
$\Alt(A)$, so fixes $A$ pointwise. If $q=1$, this proves
the second assertion of (3). Suppose $q=2$, and let $E$
be the other nontrivial orbit. Equality in (4) is forced:
\[
                         |U|=2d+1.                    \tag{5}
\]
Centralization makes $\delta$ preserve $E$. If it were
nontrivial there, it would be fixed-point-free, so $E$
would consist of $|E|/2$ pairs of $\delta$.

Partition the colours of $J$ into $L$, those with an
ambient-eligible edge in $E$, and $B$, the remaining colours.
Choose one such edge for each colour of $L$ and project it
to the pair set. Its endpoints are distinct pairs: the actual
element $\tau\in G_i$ induces $\delta$ on $E$. These
projected edges form a forest. An alternate path avoiding
colour $i$ would lift using the other selected generators
and $\tau$, all belonging to $G_i$, and would connect the
eligible edge's endpoints inside one $G_i$-orbit.
Thus $|L|\le |E|/2-1$.

For every colour of $B$, both eligible edges are confined
to $A'$. The complete selected-localization theorem for
arbitrary subsets of the interval of $J$ applies to this
single actual $J$-orbit, giving
$2|B|\le |A'|-1=|A|-6$. Therefore (5) gives
\[
 d-3=|L|+|B|
 \le\frac{|E|}{2}-1+\frac{|A|-6}{2}
 =\frac{2d+1-8}{2}=d-\frac72,
\]
a contradiction. So $\delta$ fixes $E$ as well, proving (3).
This is the acquired exterior-deficit mechanism with its
actual omitted generators and localized orbit specified.
\end{proof}

\begin{theorem}
The two opposite actual six-point heads cannot occur in
the critical configuration.
\end{theorem}
\begin{proof}
We retain the actual witness construction of the odd separator
inside $C$, rather than infer an additional property from its
numerical statement. First suppose it uses two witnesses
intransitive under $K_C$.
All their points are fixed by $\delta$. If their endpoints
are equal to $s$, the seven left head points and the two
whole orbits supply at least $2s+1$ $\delta$-fixed points
in the local fixed set. If their endpoints are unequal and
the smaller endpoint $s$ is at least $j+1$, the whole shorter
orbit and the proved truncated longer orbit supply at least
$2s$ such points, again including the seven head points.

For the remaining crossing-endpoint case the shorter orbit
ends at $j$. The separator at $s=j+1$ uses its whole orbit
and the two left endpoint orbits in the longer witness.
Their respective sizes are at least $j-3$ and $j-3$,
so with the seven head points they give
$2j+1=2s-1$ $\delta$-fixed points.
This includes $j=5,s=6$, where both sizes are exactly two.
If that exceptional case instead chooses the longer endpoint
$s=y$, its proof gives at least $2y+2$ points fixed by both
$\delta$ and the later tail, including the actual transported
orbit supplied there. These are all the intransitive-witness
cases. Lemma~2 transfers each displayed set to $F$.

Thus at least $2s-1$ points of $F$ are $\delta$-fixed.
By (3), $\delta$ fixes $U$ and at most three points of $F$
remain available. Since $\delta$ is an odd involution,
it is a literal transposition. This is excluded by the
complete exterior-pair obstruction inside $C$: it has odd
degree $2p+1$, its actual natural alternating tail and left
head are unchanged, and every required later colour is paired
in its full local ambient sense, as checked in Lemma~2.

It remains that the separator uses a $K_C$-transitive witness
$O$. Its actual transport gives a disjoint $K_C$-orbit $T$
in a $K_C$-intransitive witness. The complete proof gives
\[
 s\ge6,\qquad |O|=|T|=N\ge s-4,
 \qquad D_L\cup\{\beta,\beta^\sigma\},\ O,\ T\subseteq F,
\]
with these sets disjoint. The eight displayed head points
and $T$ are fixed pointwise by $\delta$.
If $\delta$ preserves $O$, its restriction there equals
the actual $\tau\in C_j$. It preserves the unique ambient
$C_j$-eligible witness edge and fixes its endpoints individually.
Centralization of the transitive marked $K_C$-action then
makes it fix all of $O$. At most
$|F|-(8+2N)\le2$ points remain, so oddness again makes
$\delta$ the already excluded literal transposition.

Otherwise $O^\delta$ is a distinct actual $K_C$-orbit in $F$,
disjoint from $O,T$ and the eight fixed points. It must fit
in at most two remaining points. Thus $N=2$, $s=6$, and
$F$ consists exactly of the eight points and the three
two-point orbits. The involution exchanges $O,O^\delta$
by two transpositions and fixes everything else in $F$;
it also fixes $U$. It is even, a final contradiction.
Every witness route is therefore impossible, proving the theorem.
\end{proof}

\paragraph{Attribution and remaining scope.}
The source is the opposite-end branch of Proposition~5.32 in
Cameron--Fernandes--Leemans--Mixer, \emph{Highest rank of a polytope
for $A_n$}, arXiv:1605.09173v1. The complete critical heads,
actual prefix kernels, paired intervals, odd separator proofs and
exterior-pair obstruction remain acquired. The present argument
proves the actual fixed-point transfer, combines one separator
with the reversed prefix orbit deficit, and uses the odd exterior
involution on the resulting fixed set. It requires neither an
unproved ordering of two separators nor a new induction or finite
classification. The terminal singleton attachment and full L044
source and induction assembly remain separate. No whole-original,
question, numerical or formal credit is claimed here.
\end{document}
